\chapter{Vulnerability 5.06b: Wrong Number of Arguments}

\section{Executive Summary}

As discussed in Chapter 6, vulnerability 5.06b involves a function called with more arguments than its signature declares. This violates ISO-IEC Rule 5.06 and represents an ABI mismatch where registers R2 and R3 contain uninitialized or misaligned values. This section documents the bytecode-level evidence and runtime confirmation.

\begin{tcolorbox}[colback=blue!5!white,colframe=blue!75!black,title=Key Finding,fonttitle=\bfseries]
The eBPF verifier permits function calls with too many arguments, allowing extra register values to be passed even when the function does not expect them.
\end{tcolorbox}

\section{Source-Level Patch}

A function is declared and called with more arguments than its signature expects. The vulnerability is introduced via a Git patch located at:

\begin{small}
\texttt{XDPs/xdp\_synproxy/patches/5\_06b\_argcomp/}
\texttt{0001-feat-5.06b-wrong-argument-count.patch}
\end{small}

Applied to \texttt{xdp\_synproxy\_kern.c}:

\vspace{0.3cm}
\noindent
\fcolorbox{black!30}{gray!15}{%
  \begin{minipage}[c]{\dimexpr\textwidth-2\fboxsep-2\fboxrule}
    {\small\ttfamily
    // Forward declaration: takes only 1 argument\\
    int helper\_function(int x);\\
    \\
    // But called with 3 arguments:\\
    result = helper\_function(arg1, arg2, arg3);
    }
  \end{minipage}%
}
\vspace{0.3cm}

\subsection{Bytecode-Level Evidence}

Using \texttt{llvm-objdump -S xdp\_synproxy\_kern.o}, the compiled bytecode shows:

\vspace{0.3cm}
\noindent
\fcolorbox{black!30}{gray!15}{%
  \begin{minipage}[c]{\dimexpr\textwidth-2\fboxsep-2\fboxrule}
    {\small\ttfamily
    ;~ \_\_builtin\_memcpy(dst, src, ETH\_ALEN);\\
    61: 71 41 05 00 00 00 00 00~ r1 = *(u8 *)(r4 + 0x5)\\
    /* ... more copy operations ... */\\
    72: 73 12 06 00 00 00 00 00~ *(u8 *)(r2 + 0x6) = r1\\
    ;~ bpf\_printk("[argcomp]: 5.6b - network copy called");\\
    73: 18 01 00 00 00 00 00 00 00 00 00 00 00 00 00 00~ r1 = 0x0 ll\\
    76: 85 00 00 00 06 00 00 00~ call 0x6
    }
  \end{minipage}%
}
\vspace{0.3cm}

The compiled bytecode shows both the network copy helper copy sequence and the argcomp printk, confirming that the vulnerable call path and its runtime signal exist before the verifier sees the program. Detailed bytecode analysis is available in the evidence logs.

\noindent
\textbf{Interpretation:} The sequence shows: (1) memory read from packet at offset 5, (2) memory write to destination at offset 6, (3) setup for \texttt{bpf\_printk} call, (4) the actual call instruction (0x6). The bytecode preserves the full sequence, confirming the compiler did not optimize away the vulnerable code path.

\section{Runtime Behavior}

\subsection{Triggering the Vulnerability}

Send SYN packets to the XDP interface:

\vspace{0.3cm}
\noindent
\fcolorbox{black!30}{gray!15}{%
  \begin{minipage}[c]{\dimexpr\textwidth-2\fboxsep-2\fboxrule}
    {\small\ttfamily
    python3 src/exploits/5\_06b\_argcomp/exploit\_506b.py <VM\_IP>
    }
  \end{minipage}%
}
\vspace{0.3cm}

\subsection{Observable Effects}

Kernel trace output:

\vspace{0.3cm}
\noindent
\fcolorbox{black!30}{gray!15}{%
  \begin{minipage}[c][0.8cm][c]{\dimexpr\textwidth-2\fboxsep-2\fboxrule}
    \texttt{[argcomp]: 5.6b - network copy called}
  \end{minipage}%
}
\vspace{0.3cm}

Repeated for each triggered packet, proving the function executes despite the argument count mismatch. The trace signal confirms that the vulnerable code path is executed at runtime.

\subsection{Why No Corruption?}

The extra arguments are simply ignored by the called function because arguments are passed via fixed registers (R1-R5), and the callee only consumes the registers it expects. However, the XDP context's register state is perturbed, and if the extra arguments contained sensitive values (e.g., kernel pointers), they might be disclosed. Additionally, subsequent code using registers R2 and R3 might rely on uninitialized or wrong values, creating a potential information disclosure pathway.

\section{Verification Evidence: Xlated Dump}

The \texttt{xlated} dump (verifier-processed bytecode) confirms the vulnerable call survives verification:

\vspace{0.3cm}
\noindent
\fcolorbox{black!30}{gray!15}{%
  \begin{minipage}[c]{\dimexpr\textwidth-2\fboxsep-2\fboxrule}
    {\small\ttfamily
    ; \_\_builtin\_memcpy(dst, src, ETH\_ALEN);\\
    61: (71) r1 = *(u8 *)(r4 +5)\\
    72: (73) *(u8 *)(r2 +6) = r1\\
    ; bpf\_printk("[argcomp]: 5.6b - network copy called");\\
    73: (18) r1 = map[id:7][0]+0\\
    76: (85) call bpf\_trace\_printk\#-115712
    }
  \end{minipage}%
}
\vspace{0.3cm}

The verifier-processed bytecode still contains the copy sequence and the printk call, so the argcomp path survives verification and remains executable at runtime.

\noindent
\textbf{Interpretation:} The verifier-processed bytecode is nearly identical to the compiled bytecode. The copy operations and the \texttt{bpf\_trace\_printk} call both survive verification unchanged. This confirms the verifier does not reject the function call despite the argument count mismatch.

\section{Classification: Runtime Trigger}

The verifier accepts mismatched argument counts, the bytecode preserves the call, and runtime executes the call. However, there is no direct memory corruption in this case and the extra arguments are ignored, not dereferenced.

\section{Exploitation Potential}

In isolation, 5.06b is a runtime trigger. However:
\begin{enumerate}[label=\arabic*.]
  \item If extra arguments contain pointers, information disclosure is possible
  \item Combined with register state inspection, could leak kernel addresses
  \item In a different code context, extra arguments might be used unsafely
\end{enumerate}

\section{PoC Code}

\texttt{src/exploits/5\_06b\_argcomp/poc\_506b\_final.py} sends SYN packets to trigger the vulnerable call path.

\textbf{Artifact Logs:}
\begin{itemize}
  \item Bytecode analysis: \texttt{5\_06b\_bytecode\_analysis.txt}
  \item Verifier output: \texttt{5\_06b\_verifier\_log.txt}
  \item Xlated dump: \texttt{5\_06b\_xlated\_dump.txt}
\end{itemize}
